\begin{thebibliography}{99}
\bibitem{bib-bib1} 
\par\noindent    Chris Barrett, Willem Heijltjes, and Guy McCusker.  The Functional Machine Calculus II: Semantics.  In Bartek Klin and Elaine Pimentel, editors, 31st EACSL Annual Conference on Computer Science Logic (CSL 2023), volume 252 of Leibniz International Proceedings in Informatics (LIPIcs), pages 10:1--10:18, Dagstuhl, Germany, 2023. Schloss Dagstuhl -- Leibniz-Zentrum für Informatik.  doi:10.4230/LIPIcs.CSL.2023.10.  

\bibitem{bib-bib2} 
\par\noindent    Christian Choffrut and Robert Mercas.  Contextual partial commutations.  Discrete Mathematics \& Theoretical Computer Science, 12, 2010.  

\bibitem{bib-bib3} 
\par\noindent    Mireille Clerbout and Michel Latteux.  Semi-commutations.  Information and computation, 73(1):59--74, 1987.  doi:10.1016/0890-5401(87)90040-X.  

\bibitem{bib-bib4} 
\par\noindent    Bob Coecke, Tobias Fritz, and Robert W Spekkens.  A mathematical theory of resources.  Information and Computation, 250:59--86, 2016.  doi:10.1016/J.IC.2016.02.008.  

\bibitem{bib-bib5} 
\par\noindent    Volker Diekert and Grzegorz Rozenberg.  The Book of Traces.  World Scientific, 1995.  doi:10.1142/2563.  

\bibitem{bib-bib6} 
\par\noindent    Matthew Earnshaw.  Languages of String Diagrams.  PhD thesis, Tallinn University of Technology, 2025.  doi:10.23658/TALTECH.2/2025.  

\bibitem{bib-bib7} 
\par\noindent    Matthew Earnshaw and Pawel Sobociński.  String Diagrammatic Trace Theory.  In MFCS, volume 272, 2023.  doi:10.4230/LIPIcs.MFCS.2023.43.  

\bibitem{bib-bib8} 
\par\noindent    Dominique Foata and Pierre Cartier.  Problèmes combinatoires de commutation et réarrangements.  Lecture notes in mathematics, 85, 1969.  

\bibitem{bib-bib9} 
\par\noindent    P. Freyd.  Algebra valued functors in general and tensor products in particular.  Colloquium Mathematicum, 14(1):89--106, 1966.  doi:10.4064/cm-14-1-89-106.  

\bibitem{bib-bib10} 
\par\noindent    Richard Garner and Ignacio López Franco.  Commutativity.  Journal of Pure and Applied Algebra, 220(5):1707--1751, 2016.  

\bibitem{bib-bib11} 
\par\noindent    Philip Hackney and Marcy Robertson.  On the category of props.  Applied Categorical Structures, 23(4):543--573, August 2015.  doi:10.1007/s10485-014-9369-4.  

\bibitem{bib-bib12} 
\par\noindent    Martin Hyland, Gordon Plotkin, and John Power.  Combining effects: Sum and tensor.  Theoretical computer science, 357(1-3):70--99, 2006.  doi:10.1016/J.TCS.2006.03.013.  

\bibitem{bib-bib13} 
\par\noindent    Alan Jeffrey.  Premonoidal categories and a graphical view of programs.  Preprint, 1997.  

\bibitem{bib-bib14} 
\par\noindent    André Joyal and Ross Street.  The geometry of tensor calculus, I.  Advances in Mathematics, 88(1):55--112, 1991.  doi:10.1016/0001-8708(91)90003-P.  

\bibitem{bib-bib15} 
\par\noindent    Paul Blain Levy, John Power, and Hayo Thielecke.  Modelling environments in call-by-value programming languages.  Information and Computation, 185(2):182--210, 2003.  doi:10.1016/S0890-5401(03)00088-9.  

\bibitem{bib-bib16} 
\par\noindent    Antoni Mazurkiewicz.  Concurrent program schemes and their interpretations.  DAIMI Report Series, 6(78), July 1977.  doi:10.7146/dpb.v6i78.7691.  

\bibitem{bib-bib17} 
\par\noindent    Antoni Mazurkiewicz.  Basic notions of trace theory.  In J. W. de Bakker, W. P. de Roever, and G. Rozenberg, editors, Linear Time, Branching Time and Partial Order in Logics and Models for Concurrency, pages 285--363, Berlin, Heidelberg, 2005. Springer Berlin Heidelberg.  

\bibitem{bib-bib18} 
\par\noindent    Paul-André Melliès.  Local states in string diagrams.  In Rewriting and Typed Lambda Calculi: Joint International Conference, pages 334--348. Springer, 2014.  doi:10.1007/978-3-319-08918-8\_23.  

\bibitem{bib-bib19} 
\par\noindent    John Power.  Premonoidal categories as categories with algebraic structure.  Theoretical Computer Science, 278(1):303--321, 2002.  Mathematical Foundations of Programming Semantics 1996.  doi:10.1016/S0304-3975(00)00340-6.  

\bibitem{bib-bib20} 
\par\noindent    John Power and Edmund Robinson.  Premonoidal categories and notions of computation.  Math. Struct. Comput. Sci., 7(5):453--468, 1997.  doi:10.1017/S0960129597002375.  

\bibitem{bib-bib21} 
\par\noindent    John Power and Hayo Thielecke.  Closed freyd- and $\kappa$-categories.  In Jirí Wiedermann, Peter van Emde Boas, and Mogens Nielsen, editors, Automata, Languages and Programming, pages 625--634, Berlin, Heidelberg, 1999. Springer Berlin Heidelberg.  

\bibitem{bib-bib22} 
\par\noindent    Mario Román.  Promonads and string diagrams for effectful categories.  In Applied Category Theory, volume 380 of Electronic Proceedings in Theoretical Computer Science, pages 344--361. Open Publishing Association, 2023.  doi:10.4204/EPTCS.380.20.  

\bibitem{bib-bib23} 
\par\noindent    Mario Román and Paweł Sobociński.  String diagrams for premonoidal categories.  Logical Methods in Computer Science, Volume 21, Issue 2, April 2025.  doi:10.46298/lmcs-21(2:9)2025.  

\bibitem{bib-bib24} 
\par\noindent    Sam Staton and Paul Blain Levy.  Universal properties of impure programming languages.  In Proceedings of the 40th Annual ACM SIGPLAN-SIGACT Symposium on Principles of Programming Languages, POPL ’13, pages 179--192, New York, NY, USA, 2013. Association for Computing Machinery.  doi:10.1145/2429069.2429091.  

\bibitem{bib-bib25} 
\par\noindent    Wieslaw Zielonka.  Notes on finite asynchronous automata.  RAIRO - Theoretical Informatics and Applications - Informatique Théorique et Applications, 21(2):99--135, 1987.  doi:10.1051/ITA/1987210200991.  

\end{thebibliography}
